%% kapa: exam=1 round=31 year=2020 period=1 form=A subject=economics subject_ko=경제학원론 seq=13 number=53
%% source: Q-Net 감정평가사 기출문제 (cst00309) artlSeq=5207219 file=extracted/1008/5207219/A형/A형/2020년도 제31회 감정평가사 1차 시험_1교시_A형.pdf
%% conversion: kordoc 4.24.51 -> tools/gen1.py (auto); review=proofread-sonnet; flags=none
\begin{question}{53}
공공재에 대한 甲과 乙의 수요함수가 각각 $P$甲 $=80-Q$, $P$乙 $=140-Q$ 이다. 이에 관한 설명으로 옳은 것을 모두 고른 것은? (단, $P$ 는 가격, $Q$ 는 수량)
\begin{framedtext}
ㄱ. $0≦Q≦80$일 때, 공공재의 사회적 한계편익곡선은 $P=220-2Q$ 이다. \par
ㄴ. $80<Q$ 일 때, 공공재의 사회적 한계편익곡선은 $P=80-Q$ 이다. \par
ㄷ. 공공재 생산의 한계비용이 50일 때, 사회적 최적 생산량은 90이다. \par
ㄹ. 공공재 생산의 한계비용이 70일 때, 사회적 최적 생산량은 70이다. \par
\end{framedtext}
\begin{choices}
  \item ㄱ, ㄴ
  \item ㄱ, ㄷ
  \item ㄴ, ㄷ
  \item ㄴ, ㄹ
  \item ㄷ, ㄹ
\end{choices}
\end{question}
\answer{2}
